% Chapter 4 - Real-Time Hand Tracking Pipeline

\section{Real-Time Hand Tracking Pipeline}\label{sec:pipeline}

The preceding chapters established the theoretical foundations on which the
system rests~(\autoref{sec:theory}) and positioned it against the competing
approaches surveyed in the state of the art~(\autoref{sec:sota}). This chapter
turns to the implementation. It describes the Python pipeline that converts a
single webcam frame into a canonical \acrshort{mano} pose --- a set of axis-angle
joint rotations independent of hand size and camera distance~(\autoref{subsec:gap})
--- and streams that pose, together with a recognised gesture and a placement
anchor, to a separate rendering process. The Unity application that consumes the
stream is the subject of the following chapter.

The presentation follows the same pipeline order as the theoretical background:
from the once-only preparation of the hand model, through detection, metric
landmark extraction, and the forward and inverse kinematic models, to the
serialisation and transport layer that closes the loop with the renderer.

\subsection{Pipeline Architecture and Module Organisation}\label{subsec:pipeline_arch}

The implementation separates into two stages with very different execution
profiles. An \emph{offline} stage runs once to convert the official \acrshort{mano}
release into a lightweight runtime format. An \emph{online} stage then processes
each webcam frame at interactive rate, and contains the entire per-frame data
path from capture to transmission.

The kinematic core is collected in a single package, \texttt{MANOkinematics},
whose modules correspond directly to the components introduced in the
theoretical background. \texttt{config} resolves the model file paths;
\texttt{prepare\_model} performs the offline conversion; \texttt{armatures}
defines the skeleton and the index bookkeeping that bridges the two landmark
conventions; \texttt{models} realises the \acrshort{mano} forward pass;
\texttt{AIK} is the analytical inverse-kinematics solver; and \texttt{smoother}
provides the adaptive filter applied to the raw landmarks. Around this package, a
set of top-level modules drives the live pipeline: \texttt{hand\_engine} wraps the
MediaPipe detector and exposes its latest result; \texttt{hand\_kinematics\_3D}
converts that result into smoothed metric landmarks and a placement anchor;
\texttt{mano\_mesh\_reconstruction} connects the solver to the forward model for a
single hand; and \texttt{mano\_unity\_streamer} serialises the solved pose and
transmits it.

In the online stage each frame flows through these modules in sequence. The
detector returns the 21 landmarks per hand; \texttt{hand\_kinematics\_3D} converts
them to a metric, filtered hand shape and a separate anchor describing where the
hand sits in the scene; \texttt{mano\_mesh\_reconstruction} feeds the shape to the
inverse-kinematics solver to recover \acrshort{mano} joint rotations; and
\texttt{mano\_unity\_streamer} packages the rotations, the gesture label, and the
anchor into a datagram and sends it to the renderer. Two diagnostic visualisers
--- a Matplotlib landmark view (\texttt{HandVisualizer3D}) and an Open3D mesh view
(\texttt{ManoVisualizer}) --- are available during development, but neither
participates in the streaming path and both are absent from the deployed
pipeline.

\begin{figure}[H]
	\centering
	\captionsetup{justification=centering}
	\includegraphics[width=0.95\linewidth, frame]{img/pipeline_architecture.png}
	\caption{Data flow of the real-time pipeline,
             from webcam capture to the Unity receiver.}
	\label{fig:pipeline-architecture}
\end{figure}
